\begin{ouverture}[suites]{8}{Chercher des inconnues et des régularités}{Partie 2 · Modéliser et mesurer}
{Quantités inconnues · Schémas en barres · Motifs évolutifs}
\objectif{Traduire les relations entre des quantités inconnues.}
\objectif{Représenter un problème par un schéma en barres pour le résoudre.}
\objectif{Repérer la structure d’un motif évolutif et prévoir une étape lointaine.}
\prerequis{Sens des quatre opérations ; programmes de calcul du CM2.}
\end{ouverture}
\pagestyle{otMain}\essentiel\fontsize{11}{15}\selectfont
\phantomsection\label{cours-8-1}
\essSection{1}{Décrire une quantité inconnue}
\begin{definition}Une quantité inconnue est un nombre que l’on ne connaît pas encore, mais qui obéit à des relations : « 3 fois plus que », « 15 de plus que », « ensemble, ils font 48 ». On peut la décrire avec des mots, un dessin ou une barre.\end{definition}
\begin{exemple}« Léo a 3 fois plus de billes que Mia ; ensemble, ils en ont 48. » On ne connaît aucun des deux nombres, mais on connaît les relations qui les lient.\end{exemple}
\essSection{2}{Représenter les relations par des barres}
\begin{savoirfaire}[Schéma en barres]On dessine une barre par quantité, avec des longueurs cohérentes : Mia, une barre ; Léo, trois barres identiques. Les quatre barres valent 48, donc une barre vaut $48:4=12$. Mia a 12 billes et Léo $3\times12=36$.\end{savoirfaire}
\label{fig-cours-8-1-cours-barres}\figureGeo{cours-barres}

\begin{propriete}Somme et écart : deux nombres ont pour somme 50 et pour écart 14. On retire l’écart : il reste $50-14=36$ pour deux barres égales. Le petit nombre vaut $36:2=18$, le grand $18+14=32$.\end{propriete}
\clearpage
\phantomsection\label{cours-8-2}
\essSection{3}{Raisonner sans connaître tous les nombres}
\begin{savoirfaire}[Remonter les étapes]« Je pense à un nombre, je le multiplie par 4, j’ajoute 7 et j’obtiens 39. » On défait les étapes dans l’ordre inverse : $39-7=32$, puis $32:4=8$. Contrôle : $8\times4+7=39$.\end{savoirfaire}
\begin{attention}On vérifie toujours que les nombres trouvés respectent toutes les conditions de l’énoncé, et pas seulement la dernière.\end{attention}
\essSection{4}{Comprendre un motif évolutif}
\begin{definition}Un motif évolutif change d’une étape à la suivante selon une règle. Pour comprendre sa structure, on repère ce qui reste fixe et ce qui s’ajoute à chaque étape.\end{definition}
\label{fig-cours-8-2-cours-motif}\figureGeo{cours-motif}

\begin{savoirfaire}[Prévoir une étape lointaine]Ici, il y a une allumette de départ, puis 3 allumettes par carré. À l’étape 10 : $1+3\times10=31$ allumettes ; à l’étape 100 : $1+3\times100=301$. On peut l’écrire : nombre d’allumettes $=1+3\times$ numéro de l’étape.\end{savoirfaire}

\clearpage\rubrique{Découverte}
\begin{blocLivre}{ex-8-1}\begin{exercice}[1]{Traduire par des barres}
Marc a 15 € de plus que Lina. Ensemble, ils ont 75 €. Compléter le schéma en barres, puis trouver la somme de chacun.
\label{fig-ex-8-1}\figureGeo{alg-1}
\end{exercice}\end{blocLivre}
\begin{blocLivre}{ex-8-2}\begin{exercice}[1]{Quatre fois plus}
Un livre et un cahier coûtent ensemble 20 €. Le livre coûte 4 fois plus cher que le cahier. Trouver le prix de chacun.
\end{exercice}\end{blocLivre}
\begin{blocLivre}{ex-8-3}\begin{exercice}[1]{Poursuivre un motif}
Voici les trois premières étapes d’un motif d’allumettes. Combien d’allumettes faut-il pour l’étape 4 ? Pour l’étape 5 ?
\label{fig-ex-8-3}\figureGeo{alg-3}
\end{exercice}\end{blocLivre}
\begin{blocLivre}{ex-8-4}\begin{exercice}[1]{Le nombre pensé}
Je pense à un nombre. Je le multiplie par 4, puis j’ajoute 7. J’obtiens 39. Quel est mon nombre ?
\end{exercice}\end{blocLivre}
\begin{blocLivre}{ex-8-5}\begin{exercice}[1]{Égalité de barres}
Sur le schéma, trois barres identiques et un morceau de 5 forment la même longueur qu’une barre de 26. Que vaut une barre identique ?
\label{fig-ex-8-5}\figureGeo{alg-5}
\end{exercice}\end{blocLivre}

\rubrique{Entraînement}
\begin{blocLivre}{ex-8-6}\begin{exercice}[2]{Somme et écart}
Deux nombres ont pour somme 50 et pour écart 14. Trouver ces deux nombres.
\end{exercice}\end{blocLivre}
\begin{blocLivre}{ex-8-7}\begin{exercice}[2]{Partage inégal}
Trois amis se partagent 96 cartes. Ali en a le double de Bo, et Cléa en a autant qu’Ali et Bo réunis. Combien de cartes chacun a-t-il ?
\end{exercice}\end{blocLivre}
\begin{blocLivre}{ex-8-8}\begin{exercice}[2]{Le motif en L}
Voici trois étapes d’un motif en forme de L. Combien de carreaux à l’étape 10 ? À quelle étape le motif compte-t-il 41 carreaux ?
\label{fig-ex-8-8}\figureGeo{alg-8}
\end{exercice}\end{blocLivre}
\begin{blocLivre}{ex-8-9}\begin{exercice}[2]{L’étape 100}
Pour le motif d’allumettes du cours (4, 7, 10…), combien faut-il d’allumettes pour l’étape 100 ? Décrire la règle avec une phrase.
\end{exercice}\end{blocLivre}
\begin{blocLivre}{ex-8-10}\begin{exercice}[2]{Les âges}
Aujourd’hui, Paul a 3 fois l’âge de sa fille et la somme de leurs âges est 48 ans. Quel âge a chacun ? Dans combien d’années Paul aura-t-il exactement le double de l’âge de sa fille ?
\end{exercice}\end{blocLivre}

\rubrique{Problèmes}
\begin{blocLivre}{ex-8-11}\begin{exercice}[3]{Les places de cinéma}
Deux adultes et trois enfants paient 41 € au cinéma. Une place adulte coûte 3 € de plus qu’une place enfant. Trouver le prix de chaque place.
\end{exercice}\end{blocLivre}
\begin{blocLivre}{ex-8-12}\begin{exercice}[3]{L’escalier de cubes}
On construit des escaliers de cubes : 1 cube à l’étape 1, 3 cubes à l’étape 2, 6 cubes à l’étape 3. Combien de cubes à l’étape 6 ? À l’étape 10 ?
\label{fig-ex-8-12}\figureGeo{alg-12}
\end{exercice}\end{blocLivre}
\begin{blocLivre}{ex-8-13}\begin{exercice}[3]{La croix qui grandit}
Une croix compte 5 carreaux à l’étape 1, 9 à l’étape 2 et 13 à l’étape 3. À quelle étape compte-t-elle 101 carreaux ?
\label{fig-ex-8-13}\figureGeo{alg-13}
\end{exercice}\end{blocLivre}
\begin{blocLivre}{ex-8-14}\begin{exercice}[3]{Les deux réservoirs}
Le réservoir A contient 3 fois plus d’eau que le réservoir B. On verse 20 L de A dans B : les deux réservoirs contiennent alors la même quantité. Combien chaque réservoir contenait-il au départ ?
\end{exercice}\end{blocLivre}
\begin{blocLivre}{ex-8-15}\begin{exercice}[3]{Vrai ou faux}
Nour affirme : « Dans le motif 4, 7, 10…, l’étape 10 contient 10 fois plus d’allumettes que l’étape 1. » A-t-elle raison ?
\end{exercice}\end{blocLivre}

\clearpage
